% statement-section-gap.tex — the prose heading tokens own their whole gap.
%
% Issue #177. A statement heading sits between two paragraphs, and a paragraph
% contributes `\parskip' when it starts in vertical mode. Both the heading and
% the text under it are paragraphs, so a heading skip passed through unadjusted
% would render half a line wider than the token that named it — on both sides.
%
% The tokens are documented as the complete gap, so this fixture measures the
% complete gap: every explicit skip between the heading's line box and its
% neighbours, which is the `\parskip' glue plus whatever the sectioning command
% asked for. TeX's interline glue is excluded because it is not a structural
% gap: two consecutive body lines receive it too.
%
% The fixture also pins three properties the heading design depends on:
%
%   * `\CDossierSection' and `\section*' must measure identically — one is a
%     wrapper over the other, and nothing may drift between the spellings;
%   * a statement heading carries no decorative rule, unlike the résumé and CV
%     section heading it shares its type role with; and
%   * each below-token must stay strictly greater than `\parskip'. `\@xsect'
%     reads a non-positive after-skip as a request for a run-in heading, so a
%     retune that let one reach `\parskip' would silently turn every statement
%     heading into a run-in one instead of failing.
\documentclass{careerdossier-statement}
\CDossierSetup{name = {Ada Lovelace}, email = {ada@example.test}}

\newbox\CDossierGapSectionBox
\newbox\CDossierGapStarBox
\newbox\CDossierGapSubsectionBox

\begin{document}

\setbox\CDossierGapSectionBox=\vbox{%
  \hsize=20em
  \linewidth=\hsize
  \noindent Preceding prose line.\par
  \CDossierSection{Research Vision}
  \noindent Following prose line.\par
}

\setbox\CDossierGapStarBox=\vbox{%
  \hsize=20em
  \linewidth=\hsize
  \noindent Preceding prose line.\par
  \section*{Research Vision}
  \noindent Following prose line.\par
}

\setbox\CDossierGapSubsectionBox=\vbox{%
  \hsize=20em
  \linewidth=\hsize
  \noindent Preceding prose line.\par
  \CDossierSubsection{A Narrower Theme}
  \noindent Following prose line.\par
}

\directlua{
  local subtypes = node.subtypes("glue")
  local sectionabove = \number\dimexpr\CDossierProseSectionAboveSkip\relax
  local sectionbelow = \number\dimexpr\CDossierProseSectionBelowSkip\relax
  local subsectionabove = \number\dimexpr\CDossierProseSubsectionAboveSkip\relax
  local subsectionbelow = \number\dimexpr\CDossierProseSubsectionBelowSkip\relax
  local parskip = \number\dimexpr\parskip\relax

  local function pt(sp)
    return tostring(math.floor(sp / 65536 * 1000 + 0.5) / 1000) .. "pt"
  end

  % Sum the explicit vertical space in each segment between line boxes, and
  % report whether a rule node was seen anywhere in the list. Interline glue is
  % skipped: it is TeX's line spacing, not a structural gap. Comments here must
  % use TeX's comment character, not Lua's: this argument is tokenised by TeX
  % first, which folds every newline into a space, so a Lua comment would
  % swallow the remainder of the program. The segment count is tracked in a
  % variable because Lua's length operator is a TeX parameter character.
  local function segments(boxnumber)
    local sums = {0}
    local last = 1
    local rule = false
    for item in node.traverse(tex.box[boxnumber].list) do
      local kind = node.type(item.id)
      if kind == "hlist" or kind == "vlist" then
        last = last + 1
        sums[last] = 0
      elseif kind == "rule" then
        rule = true
      elseif kind == "glue" then
        if subtypes[item.subtype] ~= "baselineskip" then
          sums[last] = sums[last] + item.width
        end
      elseif kind == "kern" then
        sums[last] = sums[last] + item.kern
      end
    end
    return sums, last, rule
  end

  local function check(label, boxnumber, above, below)
    local sums, boxes, rule = segments(boxnumber)
    if boxes < 4 then
      tex.error("The " .. label .. " fixture did not produce three line boxes")
      return
    end
    if rule then
      tex.error(
        "A statement " .. label .. " heading emitted a rule",
        {
          "Statement headings are prose headings and carry no decorative",
          "rule; the ruled heading belongs to the resume and CV classes."
        }
      )
    end
    local function compare(edge, measured, expected)
      if math.abs(measured - expected) > 655 then
        tex.error(
          "The " .. label .. " heading gap " .. edge .. " is not owned by its token",
          {
            "The token asks for " .. pt(expected) .. ", including the " .. pt(parskip),
            "an adjacent paragraph contributes; measured " .. pt(measured) .. "."
          }
        )
      end
    end
    compare("above", sums[2], above)
    compare("below", sums[3], below)
    return sums
  end

  local section = check("section", \number\CDossierGapSectionBox, sectionabove, sectionbelow)
  local star = check("section*", \number\CDossierGapStarBox, sectionabove, sectionbelow)
  check("subsection", \number\CDossierGapSubsectionBox, subsectionabove, subsectionbelow)

  if section and star
      and (section[2] ~= star[2] or section[3] ~= star[3]) then
    tex.error(
      "CDossierSection and section* do not render the same gaps",
      {
        "CDossierSection measured " .. pt(section[2]) .. " above and " ..
          pt(section[3]) .. " below;",
        "section* measured " .. pt(star[2]) .. " above and " ..
          pt(star[3]) .. " below."
      }
    )
  end

  for _, entry in ipairs({
    {"CDossierProseSectionBelowSkip", sectionbelow},
    {"CDossierProseSubsectionBelowSkip", subsectionbelow},
  }) do
    if entry[2] <= parskip then
      tex.error(
        entry[1] .. " is not greater than the paragraph skip",
        {
          "It is " .. pt(entry[2]) .. " and parskip is " .. pt(parskip) .. ", so the",
          "after-skip reaching LaTeX's @xsect would be zero or negative and",
          "every statement heading would silently become a run-in heading."
        }
      )
    end
  end
}

\box\CDossierGapSectionBox
\box\CDossierGapStarBox
\box\CDossierGapSubsectionBox

\end{document}
